Closure and the vector action: for $\beta = \tfrac12 c(u)c(v) - \tfrac12\ip{u}{v}$, the Clifford relations of $\Cl(V_\C)$ (\claimref{std-algebra.complex-clifford-algebra}(i); \claimref{std-algebra.vector-squares} in the real case) give $[\beta, c(w)] = \ip{v}{w}c(u) - \ip{u}{w}c(v) = c\bigl((u\wedge v)w\bigr)$, so $\operatorname{ad}_\beta$ preserves the space of vectors; for a second such $\beta'= \tfrac12 c(u')c(v') - \tfrac12\ip{u'}{v'}$, the derivation property gives $[\beta, \beta'] = \tfrac12\bigl(c(\beta u')c(v') + c(u')c(\beta v')\bigr)$, a sum of products of two vectors whose scalar parts $\tfrac12(\ip{\beta u'}{v'} + \ip{u'}{\beta v'})$ vanish by the antisymmetry of $\beta$ as a map on $V_\C$: the commutator is again a bivector. The bivectors act on $\CS$ because they are even. Injectivity and trace-freeness of $\rho$ are finite computations, performed exactly by the script: the fifteen basis bivectors have linearly independent, trace-free actions on $\CS$. (Trace-freeness also follows abstractly: closure gives $\rho(\mathfrak{so}) = \rho([\mathfrak{so},\mathfrak{so}]) \subseteq [\mathfrak{gl}(\CS), \mathfrak{gl}(\CS)] = \mathfrak{sl}(\CS)$; the script makes it unconditional.) Since $\dim \mathfrak{so}(2,4)_\C = 15 = \dim\mathfrak{sl}(\CS)$, $\rho$ is an isomorphism. Group level: the image of $\mathrm{Spin}(V_\C)$ is the subgroup of $GL(\CS)$ generated by the exponentials $e^{\rho(\beta)}$, and since $\rho$ is onto $\mathfrak{sl}(\CS)$ and $\mathrm{Spin}(V_\C)$ contains $e^{\beta}$ for every bivector $\beta$, it contains every $e^{tE_{ij}} = 1 + tE_{ij}$ ($i \neq j$); these elementary transvections generate $SL(4,\C)$ (row reduction: any matrix of determinant one is reduced to the identity by them), so the image is all of $SL(\CS)$. Transitivity of $SL(4,\C)$ on nonzero vectors is elementary (extend $u/\lambda$ to a basis of determinant one).
